%%%%%%%%%%%%%%%%%%%%%%%%%%%%%%%%%%%%%%%%%
% Freeman Curriculum Vitae
% German ATS Test Version (Fictional Data)
%%%%%%%%%%%%%%%%%%%%%%%%%%%%%%%%%%%%%%%%%

\documentclass[
10pt,
]{FreemanCV}

\usepackage{tikz}
\usepackage{graphicx}

\columnratio{0.55,0.45}

\begin{document}

%----------------------------------------------------------------------------------------
% NAME AND TITLE
%----------------------------------------------------------------------------------------

\parbox[][0.13\textheight][c]{\linewidth}{
	\centering
	\begin{minipage}[c]{0.72\linewidth}
		\centering
		{\sffamily\Huge Lukas Schneider}\\[6pt]
		{\cursivefont\Huge\textcolor{headings}{Lebenslauf}}
	\end{minipage}
}

%----------------------------------------------------------------------------------------
% CONTACT DETAILS
%----------------------------------------------------------------------------------------

\parbox[top][0.11\textheight][c]{\linewidth}{
	\colorbox{shade}{
		\begin{supertabular}{@{\hspace{3pt}} p{0.05\linewidth} | p{0.775\linewidth}}

		\raisebox{-1pt}{\faHome} &
		Musterstraße 42, 40210 Düsseldorf, Deutschland \\

		\raisebox{-1pt}{\faPhone} &
		+49 211 12345678 \\

		\raisebox{-1pt}{\small\faEnvelope} &
		\href{mailto:lukas.schneider@example.de}
		{lukas.schneider@example.de} \\

		\raisebox{-1pt}{\faGithub} &
		\href{https://github.com/lukas-schneider}
		{github.com/lukas-schneider}

		\end{supertabular}
	}
	\vfill
}


%----------------------------------------------------------------------------------------
% EDUCATION
%----------------------------------------------------------------------------------------

\section{Ausbildung}

\begin{supertabular}{r l}

\qualificationentry
{2019 -- 2023}
{Bachelor Informatik}
{Tutor für die Lehrveranstaltungen "Betriebssysteme" und "Signalverarbeitung"}
{Fachbereich Informatik}
{Technische Universität Düsseldorf}

\end{supertabular}


%----------------------------------------------------------------------------------------
% WORK EXPERIENCE
%----------------------------------------------------------------------------------------

\section{Berufserfahrung}


\jobentry
{Juli 2024 -- Heute}
{VZ}
{RheinTech Hosting GmbH, Düsseldorf}
{System Operations Engineer}

{Administration und Überwachung von Linux-Servern in produktiven Hosting-Umgebungen.
Erstellung technischer Berichte zur Systemverfügbarkeit.
Installation, Konfiguration und Wartung von Serverdiensten.
Verwaltung virtueller Maschinen.
Zusammenarbeit mit dem Technical Operations Center und Support-Teams zur Lösung kritischer Incidents.
Durchführung strukturierter Fehleranalysen unter Zeitdruck.}


\jobentry
{August 2023 -- Juli 2024}
{VZ}
{RheinTech Hosting GmbH, Düsseldorf}
{Network Support Engineer}

{Analyse und Behebung von Netzwerkproblemen.
Durchführung von DNS-Prüfungen und Infrastrukturdiagnosen.
Administration von cPanel- und DirectAdmin-Systemen.
Linux-Systemadministration und technische Kundenunterstützung.
Bearbeitung von Supportfällen unter Einhaltung definierter SLA-Ziele.}


\jobentry
{Januar 2021 -- Januar 2022}
{VZ}
{Language Center Düsseldorf}
{Englischlehrer}

{Unterricht von Englisch für Jugendliche und Erwachsene.
Entwicklung interaktiver Lernmaterialien.
Planung und Durchführung von Sprachkursen.}


\jobentry
{September 2022}
{TZ}
{SoftSolutions GmbH}
{Softwareentwickler Praktikant}

{Entwicklung automatisierter Webtests mit Selenium und Python.
Analyse von Benutzeroberflächen und Unterstützung bei Qualitätssicherungsprozessen.}
%----------------------------------------------------------------------------------------
% COMPUTER SKILLS
%----------------------------------------------------------------------------------------

\section{IT-Kenntnisse}

\begin{supertabular}{r l}

\tableentry{Experte}{Linux-Systemadministration}{spaceafter}

\tableentry{Experte}{Git Versionsverwaltung}{spaceafter}

\tableentry{Experte}{Rust Programmierung}{spaceafter}

\tableentry{Experte}{JavaScript / TypeScript}{spaceafter}

\tableentry{Fortgeschritten}{Python}{spaceafter}

\tableentry{Fortgeschritten}{Docker und Docker Compose}{spaceafter}

\tableentry{Fortgeschritten}{Django / Flask Webentwicklung}{spaceafter}

\tableentry{Fortgeschritten}{LaTeX und wissenschaftliches Publizieren}{spaceafter}

\end{supertabular}


%----------------------------------------------------------------------------------------
% PROJECTS
%----------------------------------------------------------------------------------------

\section{Projekte}


\subsection{\href{https://github.com/lukas-schneider/wifi-scan}
{ESP32 WLAN-Scanner -- Rust}}

Entwicklung eines leichtgewichtigen Webservers für den ESP32-Mikrocontroller.
Das System erkennt verfügbare WLAN-Netzwerke und stellt die Ergebnisse über eine
Weboberfläche dar.
Implementierung mit Rust für sichere Speicherverwaltung und effiziente Ausführung
auf eingebetteten Systemen.


\subsection{\href{https://github.com/lukas-schneider/rssh}
{SSH-Verwaltungstool -- Systemprogrammierung mit Rust}}

Entwicklung eines sicheren Kommandozeilenwerkzeugs zur Verwaltung von
SSH-Konfigurationen.

Das Tool nutzt Rust für hohe Performance und Speichersicherheit.
Es wird innerhalb der internen Support-Infrastruktur von RheinTech Hosting GmbH
eingesetzt und ersetzt eine ältere Lösung.

Durch Optimierung der Verarbeitung konnte eine Verzögerung von mehreren Sekunden
beim Zugriff auf Server entfernt werden.


\subsection{\href{https://github.com/lukas-schneider/digikala-scan}
{Web Scraper -- Node.js}}

Entwicklung eines automatisierten Datenextraktionssystems mit Node.js und
Puppeteer.

Erfassung und Verarbeitung von Produktinformationen aus webbasierten Plattformen
sowie Automatisierung wiederkehrender Datensammlungsprozesse.


\subsection{\href{https://github.com/lukas-schneider/cat}
{Cat Mimic -- C Programmierung}}

Implementierung eines Unix-\texttt{cat}-ähnlichen Kommandozeilenwerkzeugs.

Das Programm liest Dateien blockweise und ermöglicht die Verarbeitung großer
Dateien ohne unnötigen Speicherverbrauch.


\subsection{\href{https://github.com/lukas-schneider/typsy}
{typeset.live -- React und WebAssembly}}

Entwicklung einer webbasierten Plattform zur Erstellung wissenschaftlicher
Dokumente.

Implementierung eines Live-Editors mit Echtzeit-Vorschau, automatischer
Dokumentverarbeitung und direktem PDF-Export.

Die Anwendung unterstützt Studierende bei der effizienten Erstellung und
Formatierung professioneller Dokumente.


\subsection{\href{https://github.com/lukas-schneider/curltrace}
{curltrace -- eBPF / BPFtrace}}

Entwicklung eines produktionsfähigen Monitoring-Werkzeugs auf Basis von eBPF
und BPFtrace.

Das Tool wird zur Überwachung ausgehender HTTP-Anfragen auf einer Hosting-Plattform
mit mehr als 60.000 WordPress-Installationen eingesetzt.

Anwendungsbereiche:
\begin{itemize}
	\item Netzwerkdiagnose
	\item Sicherheitsanalyse
	\item Performance-Optimierung
	\item Fehleranalyse in Produktionsumgebungen
\end{itemize}


\subsection{\href{https://github.com/lukas-schneider/tourism-nextjs}
{Tourismus-Webanwendung -- Next.js}}

Entwicklung einer Full-Stack-Webanwendung mit Next.js und Prisma.

\begin{itemize}
	\item Entwicklung dynamischer Städte- und Tourangebote.
	\item Implementierung von Buchungsabläufen.
	\item Nutzung von TypeScript und Tailwind CSS für moderne Webentwicklung.
\end{itemize}


\subsection{\href{https://github.com/lukas-schneider/Dictionary-maker}
{Dictionary Maker -- Python}}

Entwicklung eines Programms zur automatischen Erstellung von Glossaren aus
Vokabellisten.

Das Projekt entstand aus dem Interesse an Sprachverarbeitung und unterstützt
die strukturierte Aufbereitung von Lernmaterialien.


\subsection{\href{https://github.com/lukas-schneider/my-vocab-list}
{My Vocab List -- Java Desktop-Anwendung}}

Entwicklung einer Desktop-Anwendung zur Verwaltung persönlicher Vokabellisten.

Integration direkter Links zu Online-Wörterbüchern und Unterstützung beim
Sprachlernen.


\subsection{\href{https://github.com/lukas-schneider/retrieval-system}
{Information Retrieval System -- Python}}

Implementierung eines Dokumenten-Suchsystems auf Basis des TF-IDF-Algorithmus.

Umsetzung von Indexierung, Dokumentbewertung und effizienter Textsuche.
%----------------------------------------------------------------------------------------
% SKILLS
%----------------------------------------------------------------------------------------

\section{Fachkenntnisse}


\subsection{Zertifikate}

VITTCO Lehrzertifikat, Language Center Düsseldorf.


%----------------------------------------------------------------------------------------

\subsection{Git und Versionsverwaltung}

Erfahrung mit dem gesamten Git-Entwicklungsprozess von der Initialisierung bis
zum produktiven Deployment.

Fundierte Kenntnisse in:

\begin{itemize}
	\item Branching-Strategien
	\item Merge- und Rebase-Prozessen
	\item Konfliktlösung
	\item Git Flow und GitHub Flow
	\item Code Review und kollaborativer Softwareentwicklung
\end{itemize}

Nutzung von GitHub zur Zusammenarbeit im Team, Verwaltung von Quellcode sowie
Dokumentation von Änderungen und Fehlerbehebungen.


%----------------------------------------------------------------------------------------

\subsection{Serveradministration, Hosting und Mail-Systeme}

Praktische Erfahrung in der Administration von Linux-basierten Hosting-
Infrastrukturen mit cPanel und DirectAdmin.

Kenntnisse in:

\begin{itemize}
	\item Verwaltung von Hosting-Accounts
	\item DNS-Konfiguration und Fehleranalyse
	\item Server-Ressourcenverwaltung
	\item Exim Mailserver-Konfiguration
	\item Dovecot IMAP/POP3 Administration
	\item Analyse von Zustellungsproblemen und Mail-Queues
\end{itemize}


%----------------------------------------------------------------------------------------

\subsection{Linux-Systemadministration und IT-Sicherheit}

Erfahrung mit der Absicherung produktiver Linux-Server:

\begin{itemize}
	\item Konfiguration und Optimierung von SSH-Zugängen
	\item Schlüsselbasierte Authentifizierung
	\item Absicherung von Root-Zugriffen
	\item Anpassung von SSH-Konfigurationen
	\item Einsatz von Fail2ban
	\item Verwaltung mit CSF (ConfigServer Security \& Firewall)
	\item Nutzung von LFD zur Angriffserkennung
\end{itemize}

Umsetzung von Sicherheitsmaßnahmen in einer Hosting-Umgebung mit mehreren
zehntausend Webseiten.


%----------------------------------------------------------------------------------------

\subsection{eBPF, Performance-Analyse und Monitoring}

Praktische Erfahrung mit moderner Linux-Systemanalyse und Kernel-Observability.

Kenntnisse in:

\begin{itemize}
	\item Entwicklung von eBPF-basierten Monitoring-Werkzeugen
	\item Analyse von Systemaufrufen und Netzwerkaktivitäten
	\item Erstellung und Interpretation von Flame Graphs
	\item CPU- und Off-CPU-Profiling
	\item PHP-Anwendungsanalyse mit Xdebug
\end{itemize}

Einsatz dieser Technologien zur Identifikation von Performance-Problemen in
WordPress-Anwendungen und produktiven Websystemen.


%----------------------------------------------------------------------------------------

\subsection{WordPress Administration und Fehleranalyse}

Umfangreiche Erfahrung bei der Analyse und Lösung komplexer WordPress-Probleme
in einer Hosting-Umgebung mit über 60.000 Webseiten.

Schwerpunkte:

\begin{itemize}
	\item Plugin- und Theme-Konflikte
	\item Datenbankoptimierung
	\item PHP-Fehleranalyse
	\item Malware-Erkennung
	\item Analyse von .htaccess-Konfigurationen
	\item Performance-Optimierung
	\item Nutzung von WP-CLI
\end{itemize}


%----------------------------------------------------------------------------------------

\subsection{Incident Management und Teamarbeit}

Erfahrung in dynamischen Produktionsumgebungen mit hohen Anforderungen an
Verfügbarkeit und schnelle Problemlösung.

Zusammenarbeit mit verschiedenen Teams:

\begin{itemize}
	\item Technical Operations Center (T.O.C.)
	\item Support-Team
	\item Entwicklungsteams
\end{itemize}

Schwerpunkte:

\begin{itemize}
	\item strukturierte Fehleranalyse
	\item technische Kommunikation
	\item Priorisierung kritischer Störungen
	\item Erstellung von Incident Reports
	\item kontinuierliche Verbesserung von Prozessen
\end{itemize}
%----------------------------------------------------------------------------------------

\subsection{Django Webentwicklung}

Erfahrung in der Entwicklung mehrerer Projekte mit dem Django Webframework.

Kenntnisse:

\begin{itemize}
	\item Django MTV-Architektur
	\item Django ORM
	\item Entwicklung eigener User-Modelle
	\item Template-Entwicklung
	\item Deployment von Django-Anwendungen
\end{itemize}

Bereitstellung eines Django-Projekts auf einer Cloud-Hosting-Plattform.


%----------------------------------------------------------------------------------------

\subsection{Docker und Containerisierung}

Praktische Erfahrung mit Docker und Docker Compose.

Kenntnisse in:

\begin{itemize}
	\item Erstellung und Verwaltung von Docker-Containern
	\item Containerisierung von Einzel- und Multi-Service-Anwendungen
	\item Aufbau reproduzierbarer Entwicklungsumgebungen
	\item Verwaltung containerbasierter Anwendungen
\end{itemize}


%----------------------------------------------------------------------------------------

\subsection{Information Retrieval und Datenverarbeitung}

Implementierung eines Information-Retrieval-Systems mit dem TF-IDF-Algorithmus.

Kenntnisse in:

\begin{itemize}
	\item Dokumentenindexierung
	\item Textanalyse
	\item Suchalgorithmen
	\item Ranking von Suchergebnissen
	\item Grundlagen von Data Mining
\end{itemize}


%----------------------------------------------------------------------------------------

\section{Sprachen}

\begin{supertabular}{r l}

\tableentry{Muttersprache}{Persisch}{spaceafter}

\tableentry{Fließend}{Englisch}{spaceafter}

\tableentry{Fortgeschritten}{Deutsch}{spaceafter}

\end{supertabular}


%----------------------------------------------------------------------------------------

\section{Zusätzliche Informationen}

\begin{itemize}

\item Erfahrung mit produktiven Linux-Serverumgebungen und Hosting-Infrastrukturen.

\item Sicherer Umgang mit Kommandozeilenwerkzeugen und Linux-Diagnose-Tools.

\item Erfahrung mit Netzwerkdiensten, DNS, HTTP, SMTP, IMAP und Servermonitoring.

\item Interesse an Systemprogrammierung, Cloud-Infrastruktur,
  Performance-Engineering und moderner Webentwicklung.

\end{itemize}


\end{document}
